\chapter{刑罚的体系}

% Source PDF physical page 181; see audit/source_coverage.json.
\begin{EditionOpeningPanel}
\InsertDiagram{figures/ch14_penalty.tex}
\end{EditionOpeningPanel}

\section{主刑}

[问题】老狗蛋在76周岁时，对乙女实施“人”酷刑，导致乙女死亡。对老狗蛋能否适用死刑？（酷刑案）

主刑的特点：（1）只能独立适用，不能附加适用；（2）一个罪只能适用一个主刑，不能同时适用两个主刑。

主刑的种类：管制、拘役、有期徒刑、无期徒刑、死刑。其中，管制属于限制自由刑，拘役属于短期自由刑，有期徒刑、无期徒刑属于长期自由刑，死刑属于生命刑。

\subsection*{一、管制}

第38条[管制的期限、禁止令、社区矫正］管制的期限，为三个月以上二年以下。判处管制，可以根据犯罪情况，同时禁止犯罪分子在执行期间从事特定活动，进入特定区域、场所，接触特定的人。

对判处管制的犯罪分子，依法实行社区矫正。

违反第二款规定的禁止令的，由公安机关依照《中华人民共和国治安管理处罚法》的规定处罚。

1.概念：管制是指对罪犯不予关押，但限制其一定自由，并采取社区矫正的刑罚方法。

% Source PDF physical page 182; see audit/source_coverage.json.
2.劳动问题：判处管制的罪犯可以自谋生计，不强制劳动，在劳动时与普通公民同工同酬。

3.期限问题

（1）管制的期限为3个月以上2年以下，数罪并罚时不得超过3年。（3-2-3）

（2）起算：刑期从判决执行之日起计算；判决执行前先行羁押的，羁押1日折抵刑期2日。

[注意】对于经过批准离开居住地的罪犯，经许可外出的期间，应计入执行期。

4.\CJKunderwave{对管制可以同时适用禁止令}。禁止犯罪分子在执行管制期间从事特定活动，进入特定区域、场所，接触特定的人。例如，对于未成年人因上网诱发犯罪的，可以禁止其在一定期限内进入网吧。

[总结]司法解释的要点：\footnote{2011年5月1日最高人民法院、最高人民检察院、公安部、司法部《关于对判处管制、宣告缓刑的犯罪分子适用禁止令有关问题的规定（试行）》。}

（1）附带民事赔偿义务未履行完毕，违法所得未追缴、退赔到位，或者罚金尚未足额缴纳的，\CJKunderwave{禁止从事高消费活动}。（2）禁止进入中小学校区、幼儿园园区及周边地区。

（3）未经对方同意，\CJKunderwave{禁止接触被害人}及其法定代理人、近亲属。（4）未经对方同意，\CJKunderwave{禁}\CJKunderwave{止接触证人}及其法定代理人、近亲属。（5）\CJKunderwave{禁止接触同案犯}。（6）\CJKunderwave{禁止令的执行期限，}\CJKunderwave{从管制、缓刑执行之日起计算}。（2012年第56题）

5.\CJKunderwave{对管制，必须实行社区矫正}。社区矫正是一种非监禁的刑罚执行方式，是将罪犯置于社区内，由社区矫正机构矫正其犯罪心理和行为恶习。

[总结]社区矫正与禁止令的比较：

\InsertDiagram{figures/manual_table_p182_1.tex}

\subsection*{二、拘役}

1.概念：短期剥夺罪犯人身自由的刑罚方法。

2.特点：短期剥夺人身自由，是一种短期自由刑。

3.执行机关：由公安机关就近在看守所执行。

4.劳动问题：被判处拘役的罪犯，参加劳动的，可以酌量发给报酬。

% Source PDF physical page 183; see audit/source_coverage.json.
5.期限问题：

（1）期限：拘役的期限为1个月以上6个月以下，数罪并罚时不得超过1年。(1-6-1)

（2）起算：刑期从判决执行之日起计算。

6.类似措施区分：拘役、刑事拘留、行政拘留（治安拘留）、司法拘留。

\subsection*{三、有期徒刑}

1.概念：剥夺罪犯一定期限的人身自由的刑罚方法。

2.期限。

（1）正常期限：6个月至15年。

（2）数罪并罚时的期限：\CJKunderwave{总和刑期不满35年的，最高不能超过20年}；\CJKunderwave{总和刑期在35}\CJKunderwave{年以上的，最高不能超过25年}。

（3）起算：刑期从判决执行之日起计算。

[提示】判决的两个日期：\CJKunderwave{判决生效或确定之日}\ensuremath{\rightarrow}\CJKunderwave{判决执行之日}（实际送监日）。

（4）折抵：判决执行以前先行羁押的，羁押1日折抵刑期1日。

[总结］先行羁押问题。管制、拘役、有期徒刑均存在判决执行以前先行羁押折抵刑期的问题。这里的“先行羁押”包括三种情形：

第一，因同一犯罪行为被刑事拘留、逮捕。这种羁押1日，折抵管制刑期2日，折抵拘役、有期徒刑刑期1日。

第二，因同一犯罪行为被采取指定居所的监视居住。这种监视居住1日，折抵管制刑期1日。这种监视居住2日，折抵拘役、有期徒刑刑期1日。\CJKunderwave{注意}：非指定居所的监视居住、取保候审，因为没有剥夺人身自由，不得折抵刑期。

第三，因同一违法行为被行政拘留、海关扣留等剥夺人身自由的措施。行为人因同一违法行为，先被采取行政拘留、海关扣留，后该行为被刑事审判。先前的行政拘留、海关扣留1日，折抵管制刑期2日，折抵拘役、有期徒刑刑期1日。\CJKunderwave{注意}：如果是两个独立的违法行为，则不存在折抵问题。\footnote{注意比较：（1）2008年试题卷二第9题：徐某因犯故意伤害罪，于2007年11月21日被法院判处有期徒刑1年，缓期2年执行。在缓刑考验期内，徐某伙同他人无故殴打学生傅某，致其轻微伤。当地公安局于2008年4月3日决定对徐某行政拘留15日，并于当日开始执行该行政拘留决定。行政拘留结束后，法院撤销对徐某的缓刑，决定收监执行。问题是：徐某被行政拘留的15天能否折抵刑期？[答案]不能，因为故意伤害罪是一件事，殴打傅某是另一件事。（2）2007年试题卷二第47题（行政法的题）：张某因打伤李某被公安局处以行政拘留15天的处罚，张某不服，申请行政复议。不久，受害人李某向法院提起刑事自诉，法院经审理认为张某的行为已经构成犯罪，判决拘役2个月。问题是：张某被行政拘留的15天能否折抵刑期？[答案】可以，因为这是同一个违法犯罪行为。}

\subsection*{四、无期徒刑}

1.概念：剥夺犯罪人终身自由，实行强迫劳动改造的刑罚方法。

2.特点：无期徒刑是自由刑中最严厉的刑罚方法。

3.执行机关：在监狱或其他场所执行。

4.适用特点：无期徒刑不能孤立适用，应当附加剥夺政治权利终身。

% Source PDF physical page 184; see audit/source_coverage.json.
\InsertDiagram{figures/manual_table_p184_1.tex}

\subsection*{五、死刑}

这是指剥夺罪犯生命的刑罚方法，包括立即执行与缓期两年执行。

\subsubsection*{（一）对适用对象的限制}

死刑（包括死缓）只能适用于罪行极其严重的犯罪分子。

对下列三类人不适用死刑（既不适用死刑立即执行，也不适用死缓）：

1.犯罪的时候不满18周岁的人。这是指犯罪的时候，不是指审判的时候。如果犯罪时不满18周岁，审判时已满18周岁，也不能适用死刑。（2011年第9题）

2.审判的时候怀孕的妇女

（1）对“审判的时候”应作扩大解释，包括整个羁押期间：审前羁押期间，审判期间，判决后执行期间。

（2）“怀孕的妇女”，是指在整个羁押期间怀孕过。只要在羁押期间处于怀孕状态，即便之后流产，不管是自然流产还是人工流产，都视为“怀孕的妇女”。如果在羁押前怀孕并在羁押前流产，不属于这里的“怀孕的妇女”。（2012年第11题，2025年试题）

3.审判的时候已满75周岁的人

第49条第2款审判的时候已满七十五周岁的人，不适用死刑，但以特别残忍手段致人死亡的除外。

（1）对“审判的时候”应作扩大解释，包括整个羁押期间：审前羁押期间，审判期间，判决后执行期间。

（2）例外：以特别残忍手段致人死亡的（包括故意杀人和故意伤害致人死亡），可以适用死刑。例如，本节开头“酷刑案”中，老狗蛋的手段特别残忍，对其可以适用死刑。

\subsubsection*{（二）死缓制度}

第50条判处死刑缓期执行的，在死刑缓期执行期间，如果没有故意犯罪，二年期满以后，减为无期徒刑；如果确有重大立功表现，二年期满以后，减为二十五年有期徒刑；\CJKunderwave{如果故意犯罪，情节恶劣的，报请最高人民法院核准后执行死刑；对于故意犯罪未执}\CJKunderwave{行死刑的，死刑缓期执行的期间重新计算，并报最高人民法院备案}。［本款经《刑法修正案（九）》修正]

% Source PDF physical page 185; see audit/source_coverage.json.
对被判处死刑缓期执行的累犯以及因故意杀人、强奸、抢劫、绑架、放火、爆炸、投放危险物质或者有组织的暴力性犯罪被判处死刑缓期执行的犯罪分子，人民法院根据犯罪情节等情况可以同时决定对其限制减刑。

1.没有故意犯罪，2年期满后，减为无期徒刑。

2.有重大立功，2年期满后，减为25年有期徒刑。旧法条是减为15年以上20年以下的有期徒刑。

3.如果故意犯罪，情节恶劣的，报请最高人民法院核准后执行死刑；对于故意犯罪未执行死刑的，死刑缓期执行的期间重新计算，并报最高人民法院备案。这是《刑法修正案（九）》的内容。

4.对下列两类严重的死缓犯，可以限制减刑：第一，累犯；第二，故意杀人、强奸、抢劫、绑架、放火、爆炸、投放危险物质或者有组织的暴力性犯罪。如何限制减刑，参见第十六讲“刑罚的执行和消灭”中“减刑”相关内容（第78条）。

5.期限问题。第51条：死刑缓期执行的期间，从判决确定之日起计算。死刑缓期执行减为有期徒刑的刑期，从死刑缓期执行期满之日起计算。注意：在判决之前的羁押时间不计算在2年死缓期间内。

\subsubsection*{（三）程序}

死刑除依法由最高人民法院判决的以外，都应当报请最高人民法院核准。死缓，可以由高级人民法院判决或者核准。

\section{附加刑}

[问题］狗蛋开车不慎撞死狗剩，撞伤小芳，撞坏小芳的电瓶车。狗蛋还欠铁牛10万元。法官判狗蛋构成交通肇事罪，在执行财产刑时，应遵守怎样的顺序？（交通肇事案）

附加刑的特点：既可以附加主刑适用，也可以独立适用，但是不能对附加刑再适用附加刑。

\InsertDiagram{figures/manual_table_p185_1.tex}

\subsection*{一、罚金}

第52条

[罚金数额的裁量］判处罚金，应当\CJKunderwave{根据犯罪情节决定罚金数额}。

第53条

[罚金的缴纳］罚金在判决指定的期限内一次或者分期缴纳。期满不缴纳的，强制缴纳。对于不能全部缴纳罚金的，人民法院在任何时候发现被执行人有可以执行的财产，应当随时追缴。

由于遭遇\CJKunderwave{不能抗拒的灾祸等原因}缴纳确实有困难的，\CJKunderwave{经人民法院裁定}，可以\CJKunderwave{延期缴}\CJKunderwave{纳}、酌情减少或者免除。［本条经《刑法修正案（九）》修正]

% Source PDF physical page 186; see audit/source_coverage.json.
1.概念：法院判处罪犯向国家缴纳一定数额的金钱的刑罚方法。罚金是财产刑的一种，在性质上与行政罚款不同。

2.适用方式。

（1）单科式，只单独适用罚金。这主要针对单位犯罪。例如，单位受贿罪（第387条）和单位行贿罪（第393条）。

（2）选科式，要么不适用，要么单独适用。例如，第275条规定，犯故意毁坏财物罪的，“处三年以下有期徒刑、拘役或者罚金”。在此，要么不适用罚金，要适用就只能单独适用。

（3）并科式，判处主刑同时并处罚金。例如，第326条规定，犯倒卖文物罪的，“处五年以下有期徒刑或者拘役，并处罚金”。在此，必须附加适用罚金，不能不用，也不能单独适用罚金。

（4）并科或单科式，要么附加适用，要么单独适用。例如，第216条规定，犯假冒专利罪的，“处三年以下有期徒刑或者拘役，并处或者单处罚金”。在此，要么并处罚金，要么单处罚金，但不可不处罚金。

\subsection*{二、没收财产}

第59条[没收财产的范围］没收财产是没收犯罪分子\CJKunderwave{个人所有财产}的一部或者全部。没收全部财产的，应当对\CJKunderwave{犯罪分子个人及其扶养的家属}保留必需的生活费用。

在判处没收财产的时候，不得没收属于犯罪分子家属所有或者应有的财产。

第60条［偿还债务］没收财产以前犯罪分子所负的\CJKunderwave{正当债务}，需要以没收的财产偿还的，\CJKunderwave{经债权人请求}，应当偿还。

第36条[赔偿经济损失与民事优先原则］由于犯罪行为而使被害人遭受经济损失的，对犯罪分子除依法给予刑事处罚外，并应根据情况判处赔偿经济损失。

承担民事赔偿责任的犯罪分子，同时被判处罚金，其财产不足以全部支付的，或者被判处没收财产的，应当\CJKunderwave{先承担对被害人的民事赔偿责任}。

1.概念：将罪犯个人财产的全部或部分强制、无偿收归国有，是一种严厉的财产刑。

2.没收财产的范围：没收犯罪分子个人合法所有并且没有用于犯罪的财产。

（1）注意严格区分罪犯个人所有财产和家属所有财产的界限。不得没收属于犯罪分子家属所有或者应有的财产；应有的财产，是指家庭共有中应属于家属的财产。

（2）没收全部财产时，应当给犯罪分子个人及其扶养的家属保留必要的生活费用。\CJKunderwave{注}\CJKunderwave{意}：这里的被扶养的家属，既包括未成年子女，也包括成年人。

3.执行机构：财产刑由一审法院执行。

4.执行顺序\footnote{2014年11月6日最高人民法院《关于刑事裁判涉财产部分执行的若干规定》第13条。}

第一位：赔偿被害人（人身损害赔偿、其他损失）。

第二位：民事债务（要求正当债务，经债权人请求）。

第三位：财产刑（罚金、没收财产）。

例如，本节开头“交通肇事案”中，法官判狗蛋构成交通肇事罪，在执行财产刑时，应遵守的顺序是：第一位，赔偿被害人的人身损害赔偿、物质损害赔偿。第二位，经债权

% Source PDF physical page 187; see audit/source_coverage.json.
人请求，偿还债权人民事债务。第三位，执行财产刑。

5.财产刑的并罚问题

\InsertDiagram{figures/manual_table_p187_1.tex}
\SourceFootnote{\textbf{表注①}\quad 注意1：罚金与没收财产虽然都是财产刑，但是属于种类不同的刑罚，所以应分别执行。之所以种类不同，是因为在执行方式上有差异。罚金可以分期缴纳（可以罚罪犯未来挣的钱），可以减免。但没收财产只能一次性执行（只能没收当下现有的财产），而且不能减免。

注意2：法律的变化。分别执行是2011年5月1日生效的《刑法修正案（八）》的新规定。此前的司法解释规定，没收全部财产和罚金应合并执行，前者吸收罚金。该司法解释规定在《刑法修正案（八）》出台后就不再适用了。2010年第56题D项，按照该司法解释规定是正确的，但按照新规定，答案就要修订了。2011年第57题、2012年第12题，按照新规定作答即可。}

\subsection*{三、剥夺政治权利}

1.概念：剥夺罪犯参与国家管理和政治活动的权利的刑罚方法。

2.必须附加适用的犯罪分子，有两类：一是主刑为死刑和无期徒刑的犯罪分子；二是危害国家安全的犯罪分子。注意：如果人民法院对危害国家安全的犯罪分子独立适用剥夺政治权利，则不能再附加剥夺政治权利。

3.起算问题

（1）独立适用时，其刑期从判决执行之日起计算并执行。

(2）附加于管制时，其刑期与管制刑期相同，同时起算（从判决执行之日起计算），同时执行。管制期满解除管制，政治权利也同时恢复。

（3）附加于拘役、有期徒刑（包括从死缓、无期徒刑改判的有期徒刑）时，其刑期从拘役、有期徒刑执行完毕之日起计算，或者从有期徒刑被假释之日起计算。\CJKunderwave{注意一}：因为拘役没有假释，所以不存在从假释之日起计算的问题。\CJKunderwave{注意二}：在拘役、有期徒刑执行期间，政治权利依然被剥夺，但不计算在剥夺政治权利刑期内。\CJKunderwave{注意三}：有期徒刑罪犯减刑时，对附加剥夺政治权利的期限可以酌减。酌减后剥夺政治权利的期限，不能少于1年。

（4）附加于死刑、无期徒刑时，其刑期是终身，从主刑执行之日起开始剥夺。

[总结]剥夺政治权利附加于拘役和管制是不同的，要注意区分。

第一，在刑期上，附加于拘役时，剥夺政治权利的期限是1年以上5年以下；附加于管制时，其期限与管制刑期相同，即3个月以上2年以下，最长3年。

第二，在起算方式上，附加于拘役时，从拘役执行完毕之日起计算；附加于管制时，从判决执行之日起计算。

\subsection*{四、驱逐出境}

1.概念：强迫犯罪的外国人离开中国国（边）境。

2.适用方式：（1）独立适用。针对犯罪情节比较轻的外国人，没有必要判处主刑的，

% Source PDF physical page 188; see audit/source_coverage.json.
可以单独判处驱逐出境。（2）附加适用。针对犯罪性质比较严重，判处了主刑的外国人，可以在主刑执行完毕后驱逐出境。

\section{非刑罚处罚措施}

[问题］狗蛋戴着黑头套，抢劫小芳，小芳反抗，狗蛋用皮鞋踢倒小芳。狗蛋构成抢劫罪。对他的黑头套、皮鞋是否可以没收？（黑头套案）

非刑罚处罚措施，是指不具有刑罚性质的处罚措施，也是承担刑事责任的一种方式。

\subsection*{一、第37条之一的措施：从业禁止}

第37条之一因\CJKunderwave{利用职业便利}实施犯罪，或者实施违背职业要求的特定义务的犯罪被判处刑罚的，人民法院可以根据犯罪情况和\CJKunderwave{预防再犯罪的需要}，禁止其\CJKunderwave{自刑罚执行完毕}\CJKunderwave{之日或者假释之日起从事相关职业}，\CJKunderwave{期限为三年至五年}。

被禁止从事相关职业的人违反人民法院依照前款规定作出的决定的，由公安机关依法给予处罚；情节严重的，依照本法第三百一十三条的规定（拒不执行判决、裁定罪）定罪处罚。

其他法律、行政法规对其从事相关职业另有禁止或者限制性规定的，\CJKunderwave{从其规定}。［《刑法修正案（九）》增设］

1.起算日期：从刑罚执行完毕之日或者假释之日起开始计算。\CJKunderwave{注意}：这里的“刑罚”是指主刑，不包括附加剥夺政治权利，也即不是从剥夺政治权利执行完毕之日起开始从业禁止。(2016年第9题）

2.第3款：其他法律有特别规定，从其规定。例如，根据律师法及相关规定，律师如果故意犯罪的，吊销执业证，终身不得从业。基于此，律师故意犯罪的，就不需要适用刑法第37条之一，适用律师法的规定即可。\footnote{根据《注册会计师法》规定（2027年1月1日生效），注册会计师犯罪（包括故意犯罪、过失犯罪），禁止从业五年；实施审计类故意犯罪，终身不得从业。}

\subsection*{二、第 64 条的措施}

第64条犯罪分子违法所得的一切财物，应当予以\CJKunderwave{追缴或者责令退赔}；对被害人的合法财产，应当及时\CJKunderwave{返还}；违禁品和供犯罪所用的本人财物，应当予以\CJKunderwave{没收}。没收的财物和罚金，一律上缴国库，不得挪用和自行处理。

\subsubsection*{（一）针对“违法所得的一切财物”}

1.违法所得的财物的范围

（1）违法所得的财物本身。例如，因赌博赢得的金钱，因受贿收受的贿赂，受雇佣杀人所得的酬金。如果是因赌博输掉的金钱，不属于违法所得。

（2）违法所得财物产生的收益。例如，将收受的贿赂用于放贷所获得的利息，用于炒股所获得的利润。

[注意］违法所得的数额以行为时为准，不受行为后的变化影响。违法所得的收益的数额以案发时为准。

% Source PDF physical page 189; see audit/source_coverage.json.
例1，甲贪污了100万元，用来炒股，赔到只剩10万元。违法所得数额仍是100万元。

例2，甲贪污了100万元，用来炒股，一个月后赚了50万元，两个月后又赔了10万元，截至案发时，共赚了30万元。犯罪所得是100万元，犯罪所得的收益是30万元。\footnote{编校注：按题面列出的两笔盈亏计算，净收益为 \(50-10=40\) 万元。原书在本例中两次写作“30万元”。}

例3，甲挪用公款100万元，炒股获利500万元。挪用的100万元不属于违法所得，因为还要归还回去；获利500万元属于违法所得。

2.措施

（1）追缴。针对尚存的违法所得的财物，进行追缴。追缴后，属于被害人的，返还被害人，其余上缴国库。\CJKunderwave{注意}：这里的追缴不同于作为附加刑的“没收财产”。后者没收的对象是犯罪人合法所有并且没有用于犯罪的财产。（2010年第56题）

（2）责令退赔。违法所得的财物已被用掉、毁掉或挥霍，无法追缴的，责令退赔。

例如，甲挪用公款300万元，炒股获利500万元，用来购房。这500万元已经不复存在，无法追缴，只能责令退赔。（2016年第8题）

[注意]共同犯罪的违法所得。例如，甲乙共同入户盗窃，甲窃得1万元，乙窃得5万元，甲乙的盗窃所得为6万元，二人均对6万元负责。但在追缴时，不能向二人各追缴6万元。（2016年第8题）

\subsubsection*{（二）没收“供犯罪所用的本人财物}

1.包括犯罪工具。例如，杀人用的刀具、走私用的船只、非法行医用的医疗器材等。

2.包括组成犯罪行为之物。例如，走私的货物、行贿的财物、聚众赌博的赌资（赌注、筹码）等。赌博输掉的金钱，不能被视为供犯罪所用的本人财物。

3.具有用于该类犯罪的通常性。

例1，甲用价值100万元的SUV故意撞毁他人价值1万元的茅草屋，该SUV不应被没收。

例2，乙用自己的房屋容留卖淫，该房屋不应被没收。

例3，本节开头“黑头套案”中，狗蛋的黑头套应该被没收，皮鞋不应被没收。

例4，丁醉酒驾车，该车不应被没收。但若丁经常醉酒驾车，多次触犯危险驾驶罪，则该车可被没收。

[注意］上述追缴或责令退赔违法所得的财物、没收违禁品和供犯罪所用的本人财物，与附加刑中的“没收财产”是有区别的。后者是一种刑罚措施，没收的对象是犯罪人个人合法所有并且没有用于犯罪的财产。
